% Einheit 2 von 5 – Dezimalzahlen addieren und subtrahieren (bank/bruchrechnung/e2.jsonl, zusammenbau v0.1)
%% TODO Zweigzeile Teil 1: was hier gelernt wird (Satz in Schülersprache aus dem Einheitstitel) – Einheit 2
\einheitenkopf[e2]{Einheit 2 von 5 · Dezimalzahlen addieren und subtrahieren}
\zweigzeile{ab Kl. 5, je nach Buch bis Kl. 6 (am Gymnasium ab Kl. 6) · keine P10-Aufgabe}
\setcounter{aufgabe}{13}

%% TODO Titel: Ich-kann-Satz fehlt in der Bank (Platzhalter: Kettenname) – Nr. 14
%% TODO Anweisung: ein Satz über der ersten Teilaufgabe fehlt in der Bank – Nr. 14
\begin{aufgabe}{Dezimal addieren/subtrahieren}
\swa{Welche Zahl steht in der Stellenwerttafel? \leerfeld}{\sachtabelle{ccc}{Einer & Zehntel & Hundertstel}{4 & 3 & 6}}
%% TODO Darstellung für schwach fehlt in der Bank (bruchrechnung-e2-k1-s1-v1; Abschnitt „Für schwache Schüler“)
\swz{$3{,}4 + 2{,}5$}{}
%% TODO Darstellung für schwach fehlt in der Bank (bruchrechnung-e2-k1-s1-v2; Abschnitt „Für schwache Schüler“)
\swz{$6{,}8 - 4{,}3$}{}
%% TODO Darstellung für schwach fehlt in der Bank (bruchrechnung-e2-k1-s1-v3; Abschnitt „Für schwache Schüler“)
\swz{$1{,}25 + 3{,}42$}{}
%% TODO Darstellung für schwach fehlt in der Bank (bruchrechnung-e2-k1-s1-v4; Abschnitt „Für schwache Schüler“)
\swz{$7{,}86 - 2{,}51$}{}
%% TODO Darstellung für schwach fehlt in der Bank (bruchrechnung-e2-k1-s1-v5; Abschnitt „Für schwache Schüler“)
\swz{$12{,}3 + 5{,}6$}{}
\swfrage{Was bleibt gleich, was ändert sich?}
%% TODO Darstellung für schwach fehlt in der Bank (bruchrechnung-e2-k1-s2-v1; Abschnitt „Für schwache Schüler“)
\swz{$5{,}6 + 1{,}23$}{}
%% TODO Darstellung für schwach fehlt in der Bank (bruchrechnung-e2-k1-s3-v1; Abschnitt „Für schwache Schüler“)
\swz{$2{,}68 + 3{,}57$}{}
%% TODO Darstellung für schwach fehlt in der Bank (bruchrechnung-e2-k1-s4-v1; Abschnitt „Für schwache Schüler“)
\swz{$3{,}45\text{ €} + 1{,}80\text{ €}$}{}
\swz[3]{Familie Aydin hat im letzten Jahr $1\,847{,}60$ € für Heizöl bezahlt, in diesem Jahr $2\,016{,}35$ €. Zeige, dass sie in diesem Jahr $168{,}75$ € mehr bezahlt hat. \hfill (P10 2014 OS)}{}
\end{aufgabe}

%% TODO Titel: Ich-kann-Satz fehlt in der Bank (Platzhalter: Kettenname) – Nr. 15
%% TODO Anweisung: ein Satz über der ersten Teilaufgabe fehlt in der Bank – Nr. 15
\begin{aufgabe}{ganze Zahl plus Dezimalzahl}
%% TODO Darstellung für schwach fehlt in der Bank (bruchrechnung-e2-k2-s1-v1; Abschnitt „Für schwache Schüler“)
\swz{$4 + 2{,}35$}{}
\end{aufgabe}

%% TODO Titel: Ich-kann-Satz fehlt in der Bank (Platzhalter: Kettenname) – Nr. 16
%% TODO Anweisung: ein Satz über der ersten Teilaufgabe fehlt in der Bank – Nr. 16
\begin{aufgabe}{Überschlag vorab}
%% TODO Darstellung für schwach fehlt in der Bank (bruchrechnung-e2-k3-s1-v1; Abschnitt „Für schwache Schüler“)
\swz{Erst überschlagen, dann rechnen: $4{,}9 + 3{,}2$}{}
\end{aufgabe}

%% TODO Titel: Ich-kann-Satz fehlt in der Bank (Platzhalter: Kettenname) – Nr. 17
%% TODO Anweisung: ein Satz über der ersten Teilaufgabe fehlt in der Bank – Nr. 17
\begin{aufgabe}{Dezimal addieren/subtrahieren – Fehler finden, Begründen, Anwendung}
\swz[3]{Paul rechnet: $3{,}5 + 0{,}42 = 0{,}77$. Finde den Fehler.}{}
\swfrage{Warum schreibt man beim Addieren Komma unter Komma?}
\swz[3]{Auf dem Kassenzettel stehen Brot 2,79 €, Käse 3,45 € und Milch 1,19 €. Du zahlst mit einem 10-€-Schein. Wie viel Rückgeld bekommst du?}{}
\end{aufgabe}

\uebersichtskasten{Komma unter Komma schreiben, fehlende Stellen mit Nullen auffüllen, dann wie mit ganzen Zahlen rechnen. \\ 4,7 + 0,85 = 4,70 + 0,85 = 5,55 \quad 6,2 − 1,35 = 6,20 − 1,35 = 4,85}
